% Options for packages loaded elsewhere
\PassOptionsToPackage{unicode}{hyperref}
\PassOptionsToPackage{hyphens}{url}
%
\documentclass[
  10pt,
]{article}
\usepackage{amsmath,amssymb}
\usepackage{iftex}
\ifPDFTeX
  \usepackage[T1]{fontenc}
  \usepackage[utf8]{inputenc}
  \usepackage{textcomp} % provide euro and other symbols
  \DeclareUnicodeCharacter{2212}{\ensuremath{-}}
\else % if luatex or xetex
  \usepackage{unicode-math} % this also loads fontspec
  \defaultfontfeatures{Scale=MatchLowercase}
  \defaultfontfeatures[\rmfamily]{Ligatures=TeX,Scale=1}
\fi
\usepackage{lmodern}
\ifPDFTeX\else
  % xetex/luatex font selection
\fi
% Use upquote if available, for straight quotes in verbatim environments
\IfFileExists{upquote.sty}{\usepackage{upquote}}{}
\IfFileExists{microtype.sty}{% use microtype if available
  \usepackage[]{microtype}
  \UseMicrotypeSet[protrusion]{basicmath} % disable protrusion for tt fonts
}{}
\makeatletter
\@ifundefined{KOMAClassName}{% if non-KOMA class
  \IfFileExists{parskip.sty}{%
    \usepackage{parskip}
  }{% else
    \setlength{\parindent}{0pt}
    \setlength{\parskip}{6pt plus 2pt minus 1pt}}
}{% if KOMA class
  \KOMAoptions{parskip=half}}
\makeatother
\usepackage{xcolor}
\usepackage[margin=1in]{geometry}
\usepackage{color}
\usepackage{fancyvrb}
\newcommand{\VerbBar}{|}
\newcommand{\VERB}{\Verb[commandchars=\\\{\}]}
\DefineVerbatimEnvironment{Highlighting}{Verbatim}{commandchars=\\\{\}}
% Add ',fontsize=\small' for more characters per line
\newenvironment{Shaded}{}{}
\newcommand{\AlertTok}[1]{\textcolor[rgb]{1.00,0.00,0.00}{\textbf{#1}}}
\newcommand{\AnnotationTok}[1]{\textcolor[rgb]{0.38,0.63,0.69}{\textbf{\textit{#1}}}}
\newcommand{\AttributeTok}[1]{\textcolor[rgb]{0.49,0.56,0.16}{#1}}
\newcommand{\BaseNTok}[1]{\textcolor[rgb]{0.25,0.63,0.44}{#1}}
\newcommand{\BuiltInTok}[1]{\textcolor[rgb]{0.00,0.50,0.00}{#1}}
\newcommand{\CharTok}[1]{\textcolor[rgb]{0.25,0.44,0.63}{#1}}
\newcommand{\CommentTok}[1]{\textcolor[rgb]{0.38,0.63,0.69}{\textit{#1}}}
\newcommand{\CommentVarTok}[1]{\textcolor[rgb]{0.38,0.63,0.69}{\textbf{\textit{#1}}}}
\newcommand{\ConstantTok}[1]{\textcolor[rgb]{0.53,0.00,0.00}{#1}}
\newcommand{\ControlFlowTok}[1]{\textcolor[rgb]{0.00,0.44,0.13}{\textbf{#1}}}
\newcommand{\DataTypeTok}[1]{\textcolor[rgb]{0.56,0.13,0.00}{#1}}
\newcommand{\DecValTok}[1]{\textcolor[rgb]{0.25,0.63,0.44}{#1}}
\newcommand{\DocumentationTok}[1]{\textcolor[rgb]{0.73,0.13,0.13}{\textit{#1}}}
\newcommand{\ErrorTok}[1]{\textcolor[rgb]{1.00,0.00,0.00}{\textbf{#1}}}
\newcommand{\ExtensionTok}[1]{#1}
\newcommand{\FloatTok}[1]{\textcolor[rgb]{0.25,0.63,0.44}{#1}}
\newcommand{\FunctionTok}[1]{\textcolor[rgb]{0.02,0.16,0.49}{#1}}
\newcommand{\ImportTok}[1]{\textcolor[rgb]{0.00,0.50,0.00}{\textbf{#1}}}
\newcommand{\InformationTok}[1]{\textcolor[rgb]{0.38,0.63,0.69}{\textbf{\textit{#1}}}}
\newcommand{\KeywordTok}[1]{\textcolor[rgb]{0.00,0.44,0.13}{\textbf{#1}}}
\newcommand{\NormalTok}[1]{#1}
\newcommand{\OperatorTok}[1]{\textcolor[rgb]{0.40,0.40,0.40}{#1}}
\newcommand{\OtherTok}[1]{\textcolor[rgb]{0.00,0.44,0.13}{#1}}
\newcommand{\PreprocessorTok}[1]{\textcolor[rgb]{0.74,0.48,0.00}{#1}}
\newcommand{\RegionMarkerTok}[1]{#1}
\newcommand{\SpecialCharTok}[1]{\textcolor[rgb]{0.25,0.44,0.63}{#1}}
\newcommand{\SpecialStringTok}[1]{\textcolor[rgb]{0.73,0.40,0.53}{#1}}
\newcommand{\StringTok}[1]{\textcolor[rgb]{0.25,0.44,0.63}{#1}}
\newcommand{\VariableTok}[1]{\textcolor[rgb]{0.10,0.09,0.49}{#1}}
\newcommand{\VerbatimStringTok}[1]{\textcolor[rgb]{0.25,0.44,0.63}{#1}}
\newcommand{\WarningTok}[1]{\textcolor[rgb]{0.38,0.63,0.69}{\textbf{\textit{#1}}}}
\usepackage{longtable,booktabs,array}
\usepackage{calc} % for calculating minipage widths
% Correct order of tables after \paragraph or \subparagraph
\usepackage{etoolbox}
\makeatletter
\patchcmd\longtable{\par}{\if@noskipsec\mbox{}\fi\par}{}{}
\makeatother
% Allow footnotes in longtable head/foot
\IfFileExists{footnotehyper.sty}{\usepackage{footnotehyper}}{\usepackage{footnote}}
\makesavenoteenv{longtable}
\setlength{\emergencystretch}{3em} % prevent overfull lines
\providecommand{\tightlist}{%
  \setlength{\itemsep}{0pt}\setlength{\parskip}{0pt}}
\setcounter{secnumdepth}{-\maxdimen} % remove section numbering
\ifLuaTeX
  \usepackage{selnolig}  % disable illegal ligatures
\fi
\IfFileExists{bookmark.sty}{\usepackage{bookmark}}{\usepackage{hyperref}}
\IfFileExists{xurl.sty}{\usepackage{xurl}}{} % add URL line breaks if available
\urlstyle{same}
\hypersetup{
  pdftitle={Resolve, Complete, Then Minimize: LR Parse Tables as Automata},
  pdfauthor={Anonymous submission},
  hidelinks,
  pdfcreator={LaTeX via pandoc}}

\title{Resolve, Complete, Then Minimize: LR Parse Tables as Automata}
\author{Anonymous submission}
\date{}

\begin{document}
\maketitle

\begin{abstract}

Canonical LR(1) construction preserves the context needed for correct
conflict resolution, but its tables often contain too many states to
ship. Practical generators therefore merge states early and then prevent
or repair changes caused by the merge. This paper explores the opposite
order: construct the canonical automaton, resolve conflicts there,
complete selected error entries, and minimize only the resulting
behavior.

The key representation encodes \texttt{reduce\ r\ on\ a} as an ordinary
transition labeled \texttt{a} into a state distinguished by production
\emph{r}. Shifts, gotos, and reductions then inhabit one labeled
transition system. Subset construction removes nondeterminism in
\emph{where the machine goes}, but a resulting subset may still contain
incompatible answers about \emph{what the parser should do}: item and
reduce nodes together encode a shift/reduce conflict, while two reduce
nodes encode a reduce/reduce conflict. Resolution classifies this output
nondeterminism before context is discarded. A conservative completion
then aligns reductions across same-core states, and ordinary partition
refinement computes the quotient. Completion is the classical
default-reduction transformation under a per-terminal, alignment-driven
selection policy. It may delay errors but preserves accepted-input
action traces.

Gazelle implements this pipeline using the same generic automaton module
as its lexer generator. On five tractable grammars, its resolved
canonical machine is checked action by action against GNU
Bison\textquotesingle s canonical-LR output by a bisimulation over
shifts, gotos, reductions, acceptance, and errors. After completion and
minimization, its state counts agree with Bison\textquotesingle s IELR
mode on six grammars. The largest fixture, a normalized SQLite grammar,
falls from 66,539 canonical states to 1,340; Bison IELR also has 1,340
states, while LALR has 1,319 and loses canonical distinctions. On the
evaluation machine, Gazelle completed the SQLite construction and
comparison in approximately six seconds; Bison\textquotesingle s
canonical mode did not finish within seven minutes. Equality of final
state counts remains evidence about compactness, not a claim that the
two final constructions are isomorphic or occupy the same number of
bytes. The result is a simpler route to a canonical-LR-equivalent parser
whose state count, on the bare grammars evaluated here, is the same as
IELR.

\end{abstract}

\hypertarget{1-introduction}{%
\section{1. Introduction}\label{1-introduction}}

Canonical LR(1) gives a parser generator an attractive semantic
baseline. Its states retain the complete one-token right context derived
from the grammar, so a conflict is inherent in the grammar rather than
introduced by table construction. If a generator resolves conflicts in
this automaton, the resolution is applied in precisely the contexts in
which each action arose.

The difficulty is size. The canonical automaton for the SQLite grammar
used in this paper contains 66,539 item sets. Its IELR automaton has
1,340 states. This is a question about how many states the machine has,
not how the eventual action/goto matrix is packed into bytes. Sparse-row
encodings, shared defaults, and similar storage techniques remain
valuable. Pure encodings of a fixed matrix are orthogonal and can be
applied after either machine has been constructed. Default reductions
sit at the boundary: they alter error behavior to make rows cheaper to
encode, and Gazelle repurposes that same kind of alteration to expose
state equivalence. This paper concerns the state-count objective. This
gap explains the usual construction order: build or
approximate the small machine first, then perform conflict analysis on
it. Unfortunately, merging LR(1) contexts can create new conflicts, and
it can also change how a real conflict is resolved.
Pager\textquotesingle s method restricts construction-time merges to
compatible states {[}3{]}. IELR begins from LALR and splits states whose
merges would alter canonical-LR behavior {[}4{]}. These algorithms
differ in detail, but both must reason about the interaction between
merging and resolution.

Gazelle reverses the dependency:

\begin{Shaded}
\begin{Highlighting}[]
\NormalTok{grammar}
\NormalTok{   {-}\textgreater{} item NFA}
\NormalTok{   {-}\textgreater{} subset construction}
\NormalTok{   {-}\textgreater{} canonical conflict resolution}
\NormalTok{   {-}\textgreater{} conservative completion}
\NormalTok{   {-}\textgreater{} partition refinement}
\NormalTok{   {-}\textgreater{} parse table}
\end{Highlighting}
\end{Shaded}

This order makes the semantic reference machine concrete. Resolution
happens before information is discarded. Later stages may merge states
only when the resolved transition system cannot distinguish them.

That description conceals a representational obstacle. A conventional LR
table contains transitions for shifts and gotos, but reductions are
actions attached to table cells or states. Generic DFA minimization
cannot preserve information absent from the transition relation. Gazelle
removes the obstacle by representing \texttt{reduce\ r\ on\ a} as a
transition on \texttt{a} into a per-production accepting state. A parser
action is determined by the kind of state reached: reaching an item
state shifts, while reaching a reduce state \emph{r} reduces by
production \emph{r}. If determinization reaches both kinds at once, it
has made the transition deterministic without making the
parser\textquotesingle s answer deterministic: that residual choice is
the conflict. Reduce states are distinguished in the initial partition
during minimization.

This paper makes four contributions:

\begin{enumerate}
\def\labelenumi{\arabic{enumi}.}
\tightlist
\item
  It gives an item-NFA encoding in which the complete LR table,
  including reductions, is one labeled transition system.
\item
  It identifies LR conflicts with output nondeterminism remaining in
  subset states after transition nondeterminism has been eliminated.
\item
  It shows that the classical default-reduction transformation, under a
  new alignment-driven selection policy, exposes common LALR-style
  merges to generic partition refinement while preserving valid-input
  parser behavior.
\item
  It checks the resolved canonical machine against
  Bison\textquotesingle s canonical-LR output by bisimulation, matches
  Bison\textquotesingle s IELR state counts after completion and
  minimization, and describes runtime precedence as an application of
  the representation.
\end{enumerate}

The claim is deliberately narrower than global minimality. Finding the
smallest conflict-free merge of canonical LR(1) states is NP-hard in
general {[}5{]}. Gazelle chooses one deterministic completion of the
transition system and computes the unique coarsest behavioral quotient
of that completed system. Another completion could, in principle, enable
a smaller sound quotient.

\hypertarget{2-background-and-problem-statement}{%
\section{2. Background and problem
statement}\label{2-background-and-problem-statement}}

Let a grammar production be

\begin{Shaded}
\begin{Highlighting}[]
\NormalTok{r: A {-}\textgreater{} X1 ... Xn}
\end{Highlighting}
\end{Shaded}

and let an LR(1) item be
\texttt{{[}A\ -\textgreater{}\ alpha\ .\ beta,\ a{]}}, where \texttt{a}
is the terminal that may follow this occurrence of \texttt{A}. Closure
introduces productions for a nonterminal after the dot, with lookaheads
selected from \texttt{FIRST(beta\ a)}. Goto advances the dot over one
grammar symbol. The canonical LR(1) automaton is the collection of item
sets reachable under closure and goto {[}1, 6{]}. Throughout, the
grammar is assumed reduced (every symbol reachable and productive) and
non-cyclic (no \texttt{A\ =\textgreater{}+\ A}) --- the standard
hypotheses under which no LR parser can perform an unbounded cascade of
reductions without consuming input.

At runtime an LR parser keeps a stack of automaton states and holds one
input token in hand. A shift follows a terminal transition and pushes
its target. A reduction by \texttt{A\ -\textgreater{}\ beta} removes
\texttt{\textbar{}beta\textbar{}} states, exposes a predecessor state,
and follows its transition on \texttt{A}. The lookahead token is not
consumed by a reduction; several reductions may occur before it is
finally shifted. Two standard terms recur below: a \textbf{handle} is a
right-hand-side occurrence whose replacement by its left-hand side is
one step in reversing a rightmost derivation, and a \textbf{viable
prefix} is a grammar-symbol sequence represented by the
parser\textquotesingle s state stack that can occur without passing the
right end of a handle. Lookahead determines whether a handle may be
reduced now; it does not make the occurrence a handle.

In a conventional presentation, the automaton transition relation
supplies shift and goto entries while a separate action table supplies
reductions. If state \emph{q} contains completed item
\texttt{{[}A\ -\textgreater{}\ beta\ .,\ a{]}}, cell \texttt{(q,\ a)}
receives \texttt{reduce\ A\ -\textgreater{}\ beta}. A shift/reduce
conflict occurs if a shift is also present in that cell, and a
reduce/reduce conflict occurs if two completed items install different
reductions there.

\hypertarget{21-why-merge-first-construction-is-difficult}{%
\subsection{2.1 Why merge-first construction is
difficult}\label{21-why-merge-first-construction-is-difficult}}

LALR merges canonical states with the same LR(0) core: the same
\texttt{(production,\ dot)} pairs after lookaheads are erased. Consider
the standard LR(1)-but-not- LALR grammar:

\begin{Shaded}
\begin{Highlighting}[]
\NormalTok{s = A x A | B x B | A y B | B y A;}
\NormalTok{x = T;}
\NormalTok{y = T;}
\end{Highlighting}
\end{Shaded}

After \texttt{A\ T}, one canonical state reduces \texttt{T} to
\texttt{x} on \texttt{A} and to \texttt{y} on \texttt{B}. After
\texttt{B\ T}, another state has the reverse mapping. Their LR(0) cores
are the same, but merging them unions the lookaheads and creates
reduce/reduce conflicts on both terminals. The grammar is LR(1); the
conflicts are artifacts of the approximation.

The same mechanism is more subtle when the canonical automaton already
has conflicts. Suppose a precedence declaration resolves a shift/reduce
conflict in favor of the reduction. A merge can union the conflicted
context with a same-core context in which the canonical automaton shifts
unconditionally --- the completed item is present there too, but without
that lookahead. The merged cell holds both actions, the declaration
selects the reduction, and the reduction now also fires in the context
whose canonical behavior was a plain shift. Reduce/reduce resolution by
rule order can migrate across contexts the same way. Thus ``the
conflicts were resolved'' is insufficient: the merged parser may differ
from the canonical parser under the same resolution policy.

IELR addresses this problem by tracking which lookahead contributions
make a LALR state inadequate and splitting it until resolved behavior
agrees with the canonical reference {[}4{]}. Gazelle instead constructs
that reference, resolves it directly, and asks a generic minimizer which
resolved states remain distinguishable.

\hypertarget{22-observable-behavior}{%
\subsection{2.2 Observable behavior}\label{22-observable-behavior}}

State minimization requires a precise behavioral boundary. For this paper, two
parsers have the same valid-input behavior when, for every token string
accepted by the canonical resolved parser, they perform the same ordered
sequence of shifts and reductions and therefore construct the same parse
tree. They must also reject every string rejected by that parser.

They need not report an invalid string at the same token or after the
same number of reductions. This distinction permits \textbf{spurious
reductions}: reductions taken only on paths that cannot lead to
acceptance. A spurious reduction can delay an error, but it cannot turn
an invalid string into a valid one when completion and minimization
satisfy the conditions in §4.

This equivalence is appropriate for recognition and valid-input
semantics. It does not preserve all observations a caller might make
during a failed parse. In particular, semantic actions with external
side effects may run on a doomed path before the error is reported.
Applications requiring transactional behavior on invalid input must
buffer or roll back such effects independently.

\hypertarget{3-encoding-the-complete-table-as-an-automaton}{%
\section{3. Encoding the complete table as an
automaton}\label{3-encoding-the-complete-table-as-an-automaton}}

The construction begins with an NFA whose ordinary states are LR(1)
items. For each item
\texttt{{[}A\ -\textgreater{}\ alpha\ .\ X\ beta,\ a{]}}, add a
transition labeled \texttt{X} to
\texttt{{[}A\ -\textgreater{}\ alpha\ X\ .\ beta,\ a{]}}. If \texttt{X}
is a nonterminal, add epsilon transitions to
\texttt{{[}X\ -\textgreater{}\ .\ gamma,\ b{]}} for its productions and
every \texttt{b\ in\ FIRST(beta\ a)}. This is the familiar item-NFA
formulation of canonical LR construction {[}6, 7{]}.

Gazelle adds one \textbf{reduce state} \texttt{R\_r} for every
production \emph{r} --- an accepting state in the lexer sense. For every
completed item \texttt{{[}A\ -\textgreater{}\ beta\ .,\ a{]}}, it adds

\begin{Shaded}
\begin{Highlighting}[]
\NormalTok{[A {-}\textgreater{} beta ., a] {-}{-}a{-}{-}\textgreater{} R\_r .}
\end{Highlighting}
\end{Shaded}

The reduce state is distinguished by production \emph{r}: reaching it
recognizes a completed right-hand side and emits the action ``reduce by
\emph{r}.'' This is the same role that a token-distinguished accepting
state plays in a lexer NFA. Only \texttt{R\_r} for the augmented start
production denotes acceptance of the whole parse; every other accepting
state returns control to the pushdown parser, which performs the
reduction and continues.

An equivalent reading is to extend every production by its lookahead and
let the dot advance one more position:

\begin{Shaded}
\begin{Highlighting}[]
\NormalTok{A {-}\textgreater{} beta . a  {-}{-}a{-}{-}\textgreater{}  A {-}\textgreater{} beta a .}
\end{Highlighting}
\end{Shaded}

All NFA edges are therefore either epsilon closure edges or ordinary dot
advances. The extra completed position is shared by all lookaheads of a
production because the runtime needs only the production identity after
the edge has been followed.

Subset construction now produces states containing both item nodes and
reduce nodes. Its transition relation has a direct runtime
interpretation:

\begin{itemize}
\tightlist
\item
  an edge on a nonterminal whose target contains items is a goto;
\item
  an edge on a terminal whose target contains items is a shift;
\item
  an edge on a terminal whose target is \texttt{R\_r} is reduction
  \emph{r}.
\end{itemize}

For a conflict-free LR(1) grammar, every reachable target used by the
parser is either an item set or one reduce node. If a subset contains
incompatible kinds of answers, determinization has succeeded as an
automaton construction but has not settled the parser action. That case
is §4.1. Otherwise the parser loop is small:

\begin{Shaded}
\begin{Highlighting}[]
\NormalTok{push(token a):}
\NormalTok{  loop:}
\NormalTok{    target = delta(top(stack), a)}
\NormalTok{    if target is item state q:}
\NormalTok{      push q; consume a; return}
\NormalTok{    if target is a reduce state R\_r for A {-}\textgreater{} beta:}
\NormalTok{      if r is the augmented start rule: accept}
\NormalTok{      pop |beta| states}
\NormalTok{      push delta(top(stack), A)}
\NormalTok{    otherwise:}
\NormalTok{      report an error}
\end{Highlighting}
\end{Shaded}

The stack still matters---the parser is not a finite-state
recognizer---but its table is a finite labeled transition system. The
distinction matters because generic subset construction and partition
refinement operate on the table, not on the pushdown control surrounding
it.

\hypertarget{31-correspondence-with-canonical-lr1}{%
\subsection{3.1 Correspondence with canonical
LR(1)}\label{31-correspondence-with-canonical-lr1}}

Ignoring reduce nodes, the epsilon closure and symbol transitions are
exactly those of the standard LR(1) item NFA. Subset construction
therefore yields the same reachable item sets as canonical closure/goto
construction. A completed item contributes an outgoing transition on
exactly its LR(1) lookahead, so the transition into \texttt{R\_r} exists
exactly where the conventional action table would contain
\texttt{reduce\ r}.

This gives a direct correspondence rather than a new parsing method:

\begin{itemize}
\tightlist
\item
  canonical item sets correspond to the item projection of subset
  states;
\item
  shift and goto entries correspond to item-targeting edges;
\item
  reduce entries correspond to edges into rule-distinguished reduce
  states.
\end{itemize}

The representation changes where actions live, not which actions the
canonical construction computes.

\textbf{Proposition 1 (canonical correspondence).} Projecting reduce
nodes out of each reachable item-bearing subset state yields the
canonical LR(1) item set reached by the same symbol path, and every
reachable canonical item set occurs as such a projection. Pure reduce
subsets represent action outputs and project to the empty set. For every
terminal and nonterminal, the encoded machine contains the corresponding
shift, goto, or reduce edge if and only if the conventional canonical
table contains that action. The item-bearing projection is many-to-one:
a hybrid state and a pure item state can share one item set, so the
encoded machine may carry more physical states than the canonical
automaton has item sets (§6.1 accounts for the duplicates when
counting).

The argument follows directly by induction over subset construction:
epsilon closure is LR(1) closure, symbol advance is goto, and the only
additional edge maps each completed item to the reduction already
prescribed by its lookahead.

\hypertarget{4-resolve-first-complete-then-minimize}{%
\section{4. Resolve first, complete, then
minimize}\label{4-resolve-first-complete-then-minimize}}

\hypertarget{41-the-nondeterminism-that-remains}{%
\subsection{4.1 The nondeterminism that
remains}\label{41-the-nondeterminism-that-remains}}

Subset construction is deterministic: it produces at most one outgoing
edge per symbol. But this removes only \textbf{transition
nondeterminism}. It says where the machine goes; it does not guarantee
that every NFA state collected at the destination assigns the same
meaning to arriving there.

This distinction is familiar in lexer generation. If a determinized
regex state contains accepting NFA states for both \texttt{IDENT} and
\texttt{KEYWORD}, the DFA has one transition path but two possible token
outputs. The lexer generator still needs a priority rule. Its residual
problem is not graph nondeterminism but \textbf{output nondeterminism}.

The reduce-state encoding makes LR conflicts the same phenomenon. A
conflict cannot appear as two competing DFA edges; it appears in their
union target.

If an item advances on terminal \texttt{a} while a completed item
reduces on \texttt{a}, the target of the single \texttt{a} edge contains
both advanced items and a reduce node. This is a \textbf{hybrid state}:
the deterministic destination says both ``shift'' and ``reduce
\emph{r}.'' If two completed items reduce on \texttt{a}, the target
contains two reduce nodes and says both ``reduce \emph{r1}'' and
``reduce \emph{r2}.'' The state records exactly the set of actions that
a conventional table would place in the conflicted cell.

\begin{Shaded}
\begin{Highlighting}[]
\NormalTok{shift item   {-}{-}a{-}{-}\textgreater{}  advanced item {-}{-}+}
\NormalTok{                                      +{-}{-}\textgreater{} subset \{advanced item, R\_r\}}
\NormalTok{completed    {-}{-}a{-}{-}\textgreater{}  R\_r {-}{-}{-}{-}{-}{-}{-}{-}{-}{-}{-}+}
\NormalTok{                           one edge, two parser outputs}
\end{Highlighting}
\end{Shaded}

\textbf{Figure 1.} Subset construction combines transition targets but
does not choose between their parser meanings.

Thus an LR conflict is the nondeterminism subset construction cannot and
should not erase: multiple semantic outputs survive in one deterministic
subset. The grammar, together with one token of evidence, has not
selected a unique parse action. Conflict resolution is the separate act
of choosing an output.

Gazelle first reports these states, retaining their canonical item
context for diagnostics and counterexample generation. Classification is
then driven not by a global policy but by per-terminal declarations in
the grammar. A terminal has one of five kinds:

\begin{Shaded}
\begin{Highlighting}[]
\NormalTok{plain           a conflict is an error, reported with counterexamples;}
\NormalTok{                it compiles only under a matching \textasciigrave{}expect\textasciigrave{} count}
\NormalTok{shift ELSE      shift/reduce on this terminal resolves to shift}
\NormalTok{reduce X        shift/reduce on this terminal resolves to reduce}
\NormalTok{prec OP         both actions are kept; runtime precedence decides (§5)}
\NormalTok{conflict NAME   both actions are kept; the lexer decides per token (§5)}
\end{Highlighting}
\end{Shaded}

There is no silent default. An unannotated conflict fails generation
with its counterexamples unless the grammar acknowledges the exact
conflict count (\texttt{expect\ 3\ rr;}), in which case acknowledged
reduce/reduce conflicts resolve to the earlier rule; an acknowledged
shift/reduce conflict resolves to shift. Resolution is data in the
grammar rather than an undocumented tool default: the dangling else is
the single declaration \texttt{shift\ ELSE}, naming the classical
disambiguation where the terminal is introduced. The property that
matters for this paper is timing: whatever the declarations say,
classification happens on the unmerged canonical machine. After it, each
reachable state has one runtime meaning --- or, for \texttt{prec} and
\texttt{conflict} terminals, one precisely delimited deferred question
(§5).

A hybrid state can duplicate the item behavior of a pure item state
elsewhere in the automaton. This is an artifact of putting alternatives
into target states. Once classification discards the losing reduce
nodes, ordinary minimization merges such duplicates when their outgoing
behavior agrees.

\hypertarget{42-why-minimization-alone-is-insufficient}{%
\subsection{4.2 Why minimization alone is
insufficient}\label{42-why-minimization-alone-is-insufficient}}

Canonical same-core states often differ only because their reductions
are defined on disjoint lookahead sets. Consider states \texttt{q1} and
\texttt{q2} that both reduce by \emph{r}, with \texttt{q1} defining the
reduction on \texttt{a} and \texttt{q2} defining it on \texttt{b}. Their
valid continuations may be identical even though their partial
transition rows differ:

\begin{Shaded}
\begin{Highlighting}[]
\NormalTok{       a       b}
\NormalTok{q1    R\_r     error}
\NormalTok{q2    error   R\_r}
\end{Highlighting}
\end{Shaded}

A minimizer correctly keeps these rows apart. Yet completing both gaps
with \texttt{R\_r} is harmless when \texttt{b} cannot occur at
\texttt{q1} and \texttt{a} cannot occur at \texttt{q2} on a valid parse.
Completion makes the latent equivalence visible:

\begin{Shaded}
\begin{Highlighting}[]
\NormalTok{       a       b}
\NormalTok{q1    R\_r     R\_r}
\NormalTok{q2    R\_r     R\_r}
\end{Highlighting}
\end{Shaded}

The added entries may cause reductions on invalid input. If the
reduction eventually reaches an accepting continuation, however, the
supposedly impossible lookahead would witness a valid continuation of
the original state, contradicting the condition under which the gap was
filled. Thus the new edge can postpone rejection but cannot create an
accepted sentence.

\hypertarget{43-completion-versus-table-packing}{%
\subsection{4.3 Completion versus table
packing}\label{43-completion-versus-table-packing}}

Two different compression problems are easy to conflate. \textbf{State
minimization} reduces the number of states in the parser automaton; this
is the subject of the paper. \textbf{Table packing} encodes the
action/goto matrix of an already chosen automaton in fewer bytes.
Joliat\textquotesingle s reduced-matrix work {[}11{]} belongs to the
second problem. Yacc\textquotesingle s per-state defaults {[}14{]} and
Bison\textquotesingle s default-reduction options {[}12{]} are more
closely related to this paper: they first alter the partial machine,
treating selected error cells as reductions, so the altered rows admit a
sparser representation. Their objective is still row storage rather
than state count, but their mechanism is semantic completion rather
than pure encoding.

Gazelle\textquotesingle s pre-minimization completion has the same
cell-level effect as a default reduction: an error entry is treated as a
reduction already performed elsewhere in that state. The purpose and
selection policy differ. A classical default chooses one reduction per
state so its row can be stored compactly. Gazelle chooses per terminal
across same-core states so partition refinement can identify equivalent
states. The former reduces the representation of a fixed matrix; the
latter changes the completed transition system presented to the
minimizer and thereby reduces the state set. Recognition and
accepted-input action traces cannot tell
these insertions apart: in both, an error entry is replaced by a
reduction that no successful parse consults. Failed-parse behavior can
distinguish them. The policies may run different semantic actions on
doomed paths, delay error reporting by different amounts, and change
when a feedback-sensitive lexer is called.

Under either policy, every inserted entry satisfies the same two local
conditions: (i) it replaces an error entry --- no existing action
changes --- and (ii) it extends the lookahead set of a reduction the
state already performs. Condition (ii) holds for
Gazelle\textquotesingle s rule because the donor sibling has the same
LR(0) core, so the donated reduction\textquotesingle s completed item is
already present in the receiving state; completion never imports a
reduction from elsewhere. When both techniques are modeled as completed
transition relations, classical per-state defaults are a special case
of these insertions, so the preservation argument of §4.4 covers their
semantic effect as well. This shared model does not make their
compression objectives the same. One assumption deserves note: condition (i) reads
an error entry as evidence of non-viability, which is true of the
canonical machine under the policies of §4.1. A resolution policy that
deliberately maps \emph{viable} lookaheads to error ---
yacc\textquotesingle s \texttt{\%nonassoc} --- would create error cells
that this argument does not cover, and such cells must be marked
non-fillable, exactly as §5\textquotesingle s virtual-twin guard already
does for deferred cells. This is not hypothetical: Bison identifies
\texttt{\%nonassoc}, default reductions in inconsistent states, and
parser-state merging as distinct causes of delayed or inaccurate
syntax-error behavior, and provides lookahead correction
(\texttt{parse.lac}) to address them {[}12{]}.

Concretely, Gazelle groups resolved item states by LR(0) core. Within a
group, for every terminal \emph{a}, it inspects the reduce targets
already present:

\begin{enumerate}
\def\labelenumi{\arabic{enumi}.}
\tightlist
\item
  If every state defining a reduce transition on \emph{a} agrees on the
  same target \texttt{R\_r}, add that edge to siblings where \emph{a} is
  absent.
\item
  If two states reduce different productions on \emph{a}, leave every
  gap intact.
\item
  Never replace an existing transition.
\end{enumerate}

The disagreement rule preserves the split in the LR(1)-but-not-LALR
example of §2.1. There, the same terminal selects different reductions
in the two states, so completion cannot make their rows equal.

The procedure is intentionally conservative. It does not search all
possible ways to complete error entries, nor does it claim that its
choice produces the smallest possible parser. It chooses a deterministic
completion justified by same-core context and unanimous existing
reductions.

The two objectives remain distinct, and Gazelle uses the shared
mechanism at two stages. Alignment-driven insertion runs before
minimization to reduce the number of states. Per-state defaulting runs
afterwards: each minimized state elects its most frequent reduction,
changes remaining error cells to that default on doomed paths, and omits
matching entries from the packed row. The first use is the
paper\textquotesingle s algorithmic contribution; the second is the
classical row-sparsification use of the same behavioral relaxation.

After completion, Gazelle runs iterative partition refinement. The
initial partition places all item states together and places reduce
states in classes distinguished by production. Refinement repeatedly
splits a class when two members transition, on some label, into
different current classes. At the fixed point, quotienting by the
partition produces the coarsest transition- preserving merge of the
completed machine. The implementation is Moore-style iterative
refinement despite its historical function name
\texttt{hopcroft\_minimize}.

\hypertarget{44-preservation-argument}{%
\subsection{4.4 Preservation argument}\label{44-preservation-argument}}

\textbf{Proposition 2 (resolved-behavior preservation).} Given the
completion rule of §4.3 and a deterministic classification of every
canonical conflict, the quotient parser accepts exactly the strings
accepted by the resolved canonical parser and performs the same shifts
and reductions on each such string.

There are two transformations to justify.

\textbf{Completion.} Write a parser configuration as
\texttt{(sigma,\ w)}, where \texttt{sigma} is the stack of canonical
states and \texttt{w} is the unread token string. Existing edges are
never changed. Every configuration on an accepting canonical run
therefore follows the same edge as before, and its sequence of shifts
and reductions is unchanged.

For rejection, consider the first inserted edge consulted in a run. It
is a reduction by \texttt{A\ -\textgreater{}\ beta} on lookahead
\texttt{a} from a state whose canonical cell on \texttt{a} was an error.
The source state nevertheless contains the completed LR(0) item
\texttt{A\ -\textgreater{}\ beta\ .}; same-core completion has extended
only its lookahead. Let the stack encode a viable prefix ending in
\texttt{gamma\ beta}. If taking the inserted reduction and later
accepting were possible, the accepting run would witness a rightmost
derivation in which \texttt{a} can follow the corresponding occurrence
of \texttt{A} after \texttt{beta} is reduced. Canonical LR(1) lookahead
propagation would then include \texttt{a} on that completed item and put
a reduce action in the original cell, contradicting that the cell was an
error. Thus an inserted reduction preserves the non-viability of the
held lookahead. The repetition is licensed by an invariant worth
stating: the completed LR(0) item certifies that \texttt{beta} is a
handle independently of the held lookahead, so reducing it preserves the
viable-prefix invariant even though \texttt{a} is not valid there. The
argument therefore applies verbatim at the next inserted edge. A doomed
run never reaches a shift of the held token or an accepting
configuration.

The parser may pop and reduce before discovering the error, but it
rejects without consuming the offending token: the correct-prefix
property is preserved. The argument is local to each inserted entry and
run, so choosing donors independently per terminal creates no
interaction between insertions. The reduced, non-cyclic hypothesis of §2
supplies termination: every inserted edge is a locally legal reduction,
and a no-input reduction cascade cannot be unbounded. It must end in a
state whose row has no entry for the held token.

\textbf{Quotienting.} Partition refinement merges only states with the
same classification and the same labeled transitions modulo the final
partition. Relate a completed-machine stack \texttt{q0\ ...\ qn} to the
quotient stack \texttt{{[}q0{]}\ ...\ {[}qn{]}} pointwise. A shift
preserves this relation because equivalent states have transitions on
the same label into the same target class. A reduction preserves it
because different rule states never merge, both stacks pop the same
number of entries, and equivalent exposed states have equivalent gotos
on the production\textquotesingle s left-hand side. Induction over
parser steps therefore gives the same action at each configuration and
preserves acceptance in both directions.

Together, the transformations preserve accepted strings and the action
trace on every accepted string. They may change the point at which a
rejected string fails, as allowed by §2.2.

This remains a proof sketch of the implemented criterion, not a claim
that every completion satisfying the same external semantics has been
characterized. A full proof should define viability over stack
configurations and prove the inserted-reduction lemma above directly
from the canonical item construction. A mechanized proof or an
executable equivalence checker would strengthen the result further.

\hypertarget{45-relation-to-np-hard-minimization}{%
\subsection{4.5 Relation to NP-hard
minimization}\label{45-relation-to-np-hard-minimization}}

Yang proves that minimizing LR(1) state machines is NP-hard, reducing
graph coloring to the question of which same-core states of a canonical
machine can be merged without introducing conflicts {[}5{]}. The hard
search is over merge choices on a machine that is already built: merges
interact --- combining one pair can enable or forbid combining others
--- so the minimum conflict-free merge is a global combinatorial
problem. (Each such merge can equally be read as a choice of error-entry
insertions, since merging states with disjoint reduce lookaheads is
completing each with the other\textquotesingle s entries; this
completion view is ours, not Yang\textquotesingle s.) Gazelle does not
attack that problem. It fixes the inserted entries once, by the
unanimous-reduction rule, and computes the unique coarsest quotient of
the resulting completed transition system --- polynomial and
deterministic, but possibly finer than what a different completion, or a
cleverer merge, could reach.

The word ``minimize'' in this paper always refers to this fixed
completed machine. It does not mean globally fewest states among all
sound LR tables: Yang\textquotesingle s optimum ranges over all
conflict-free merges of the canonical machine, and nothing here claims
to reach it.

\hypertarget{5-runtime-precedence-as-an-application}{%
\section{5. Runtime precedence as an
application}\label{5-runtime-precedence-as-an-application}}

Some conflicts are intentional. An expression rule such as

\begin{Shaded}
\begin{Highlighting}[]
\NormalTok{expr = expr OP expr | atom}
\end{Highlighting}
\end{Shaded}

leaves associativity and precedence to a language policy. Traditional
generators resolve each conflicted table cell statically. Gazelle can
instead defer selected shift/reduce choices to token data, allowing one
terminal to represent operators whose precedence is known only at
runtime.

All non-plain terminal declarations use the same representational
device. For a modified terminal \texttt{OP}, construction creates a
virtual symbol \texttt{OP\_reduce}. Shift transitions retain the real
symbol; completed items use the virtual symbol on edges to reduce
states:

\begin{Shaded}
\begin{Highlighting}[]
\NormalTok{item {-}{-}OP{-}{-}{-}{-}{-}{-}{-}{-}\textgreater{} item target}
\NormalTok{item {-}{-}OP\_reduce{-}\textgreater{} reduce target}
\end{Highlighting}
\end{Shaded}

Because the labels differ, subset construction and minimization see an
ordinary deterministic automaton. During table extraction,
\texttt{shift} and \texttt{reduce} declarations select their named
branch statically. The \texttt{prec} and \texttt{conflict} declarations
combine the two columns into a \texttt{shift-or-reduce} entry. For
\texttt{prec}, the pending construct carries precedence
\texttt{p\_stack} and the incoming token carries \texttt{p\_token}. A
higher \texttt{p\_token} shifts; a lower one reduces. At equal levels,
left associativity reduces and right associativity shifts. Thus, while
parsing \texttt{1\ +\ 2\ *\ 3}, the incoming \texttt{*} shifts over the
pending \texttt{+}, whereas a following \texttt{+} causes the pending
\texttt{*} expression to reduce.

This is the shunting-yard decision rule embedded in an LR parser
{[}15{]}: shift an incoming operator when it binds more tightly than the
pending operator; reduce when it binds less tightly; and on equal levels
use associativity to choose. Gazelle needs no separate operator stack or
output queue. Its LR stack carries the pending precedence, and ordinary
grammar reductions produce the result. At an individual conflicted cell,
the \texttt{conflict} form below exposes the more general choice:
instead of applying this built-in policy, the token selects the shift or
reduce branch directly.

That direct choice does not make \texttt{prec} mere shorthand in the
generated API. A \texttt{conflict} token carries one fixed
\texttt{Shift} or \texttt{Reduce} value throughout its \texttt{push}. An
incoming operator can cause several reductions before it is shifted, and
the correct shunting-yard choice must be recomputed after each one
against the newly exposed pending operator. In
\texttt{1\ +\ 2\ *\ 3\ *\ 4}, for example, the second \texttt{*} must
first reduce the pending \texttt{*}, then shift over the newly exposed
\texttt{+}; one fixed choice cannot do both. A semantic action could pop
a separate operator stack, but it cannot revise the current
token\textquotesingle s already fixed resolution. A user could reproduce
the policy with the low-level parser by resubmitting a decision after
every reduction; \texttt{prec} performs that loop inside the generated
parser, using the LR stack itself. This makes \texttt{prec} the
specialized path for the common expression-parsing case, while
\texttt{conflict} remains the escape hatch for other token-directed
choices.

Completion needs one additional guard. A state may lack
\texttt{OP\_reduce} while having a real \texttt{OP} shift. Filling the
virtual gap would not merely add a reduction on an invalid path; it
would turn a canonical unconditional shift into a deferred choice on
valid input. Gazelle therefore does not fill a virtual reduce edge into
a state that has a transition on its real twin.

\texttt{conflict} terminals use the same mechanism with a different
decision source: both branches survive to the table, and the lexer,
rather than a precedence comparison, supplies the answer per token.
Gazelle uses this explicit feedback channel for C\textquotesingle s
typedef ambiguity; the contextual decision remains in the lexer, but its
effect on parsing is represented directly in the table.

The comparison below distinguishes a \textbf{table-local resolver} from
an unrestricted callback. A table-local resolver chooses from the
current cell using the current token and bounded metadata already
maintained by the parser; it does not reread the consumed input or run a
shadow parser. Without this boundary, runtime resolution is vacuous as a
parser-generation result: a generator could emit a loop that asks
arbitrary user code for every action, and that code could implement the
canonical parser itself. Gazelle\textquotesingle s \texttt{prec}
resolver is table-local: each decision compares the incoming
token\textquotesingle s precedence with precedence already carried on
the LR stack. Gazelle\textquotesingle s \texttt{conflict} terminals are
not table-local in how their one fixed per-token answer is computed: the
lexer may consult arbitrary state. The boundary they respect is a
different one: the invariant below guarantees that a deferred cell
exists only where the canonical machine has the question. The
unrestricted answer is therefore solicited only at canonical choice
points, never merely because merging converted an unconditional
canonical action into a question. The answer may still depend on
semantic state intentionally kept outside the grammar.

Deferral makes the advantage of the construction order concrete. A
deferred cell may contain only actions available in the corresponding
canonical cell; otherwise runtime data could select an action that the
canonical machine never offered in that context. Gazelle preserves this
invariant by separating shift and reduce branches before determinization
and retaining their distinct labels through minimization; the completion
guard above enforces the same invariant at the one place completion
could otherwise violate it, by manufacturing a deferred question where
the canonical machine had an unconditional answer.

A documentation survey found one close precedent among fourteen LR
parser generators. Parglare lets productions retain dynamically filtered
actions and calls a user-supplied predicate at parse time; its manual
demonstrates runtime-defined operator precedence with the deterministic
\texttt{Parser} API. Its worked example ranks the first distinct
operator encountered lowest and each later one higher. Its callback then
shifts when the incoming operator ranks higher than the pending one and
reduces when it ranks lower or equal---the shunting-yard comparison,
with input-defined priorities and left-associative ties. That is
user-written policy rather than a built-in parser abstraction. Compared
with that general callback, Gazelle carries precedence on tokens,
performs the comparison directly in the generated parser on one stack,
and uses its real/virtual transition representation to preserve the
deferred choice through completion and minimization. The faithfulness
question also separates the two: parglare\textquotesingle s filter
selects among whatever actions its merged LALR tables retain, while
Gazelle guarantees that every deferred alternative was available in the
corresponding canonical cell --- the invariant above and the obligation
below. The distinction is observable: Appendix A.4 constructs a
five-production grammar on which canonical LR(1) shifts a token
unconditionally in one context and has a genuine deferred question on
the same token in a same-core context; parglare\textquotesingle s merged
table gives one cell to both, and either erases the question at
construction or asks it in both contexts. No table-local filter can
answer correctly on both sides. An unrestricted callback can recover by
rereading left context or maintaining a shadow parser, but doing so
moves the context reconstruction that the table omitted into user code.
Each of the other thirteen surveyed generators documents only
generation-time precedence in deterministic LR mode. GLR systems also
offer parse-time ranking of forked parses. Appendix A records the survey
and its scope.

Deferral imposes a different proof obligation on a merge-first
construction. IELR\textquotesingle s inadequacy analysis is designed to
preserve the result of static conflict resolution; analogous runtime
deferral would appear to require preserving the complete canonical
action set and whether a conflict exists in each context. We have not
implemented that extension, so this is an architectural observation
rather than a complexity or impossibility claim. With
Gazelle\textquotesingle s real/virtual branch representation in place,
preserving deferral through completion costs exactly that one guard.
Bison\textquotesingle s \texttt{\%dprec} is superficially related but
operates in GLR mode by ranking surviving parses {[}12{]}; Gazelle
instead chooses one action on a single deterministic LR stack before
either branch executes.

Gazelle\textquotesingle s regression suite compares this runtime
decision with a conventional fixed grammar hierarchy for all 1,093
operator sequences containing up to seven operands and three precedence
levels. The two parsers produce identical abstract syntax trees in every
case. This is bounded validation of the runtime mechanism, not a proof
for arbitrary grammars.

The transition system preserves both branches until runtime, and
partition refinement keeps apart exactly the states whose labeled
choices differ. A full treatment of runtime precedence---including its
interaction with lexer feedback and semantic values---is outside this
paper\textquotesingle s central claim.

\hypertarget{6-evaluation}{%
\section{6. Evaluation}\label{6-evaluation}}

Gazelle\textquotesingle s \texttt{-\/-yacc} mode emits an equivalent
Bison grammar, allowing Bison to serve as an independent implementation
reference. The evaluation uses six grammars: C11, Python, jq, SQLite,
Gazelle\textquotesingle s regular-expression grammar, and its
self-hosted meta grammar. jq is mechanically normalized from its native
Bison grammar; SQLite is normalized from the grammars-v4 ANTLR grammar.
The sources, pinned revisions, licenses, and transformations are
recorded with the fixtures. They exercise parser-table construction and
are not claims of complete Gazelle front ends.

Terminal resolution modifiers (\texttt{shift}, \texttt{reduce},
\texttt{prec}, \texttt{conflict}) and upstream static-precedence
declarations are stripped for the comparison, so both tools receive the
same bare grammar and apply the same default conflict choices. These
comparison machines are not the production Gazelle tables when a grammar
uses modifiers; §6.2 reports those counts separately.

\hypertarget{61-canonical-construction}{%
\subsection{6.1 Canonical
construction}\label{61-canonical-construction}}

The first comparison uses \texttt{bison\ -Dlr.type=canonical-lr} with
default reductions restricted to the accepting state. Bison contributes
one synthetic \texttt{\$accept} state; the table reports its count minus
that state.

\begin{longtable}[]{@{}lrr@{}}
\toprule\noalign{}
grammar & Gazelle canonical item sets & Bison canonical −
\texttt{\$accept} \\
\midrule\noalign{}
\endhead
\bottomrule\noalign{}
\endlastfoot
C11 & 2,097 & 2,097 \\
Python & 3,298 & 3,298 \\
jq & 4,779 & 4,779 \\
regex & 69 & 69 \\
meta & 61 & 61 \\
\end{longtable}

Conflict counts agree as well. More strongly, the automated regression
exports Bison\textquotesingle s resolved actions and compares the two
reachable labeled transition systems. Starting from their initial
states, every terminal and nonterminal must produce the same shift,
goto, reduction, acceptance, or error; successor pairs are added to the
relation recursively, and every item state in both machines must be
covered. The relation is intentionally not required to be a bijection
because Gazelle\textquotesingle s hybrid encoding can contain multiple
physical targets for one conventional LR state. This canonical
bisimulation passes for all five grammars in the table.

SQLite has 66,539 Gazelle canonical states. GNU Bison\textquotesingle s
canonical construction did not finish within seven minutes on the
evaluation machine, so no Bison canonical count or cross-tool
bisimulation is reported for that fixture. It is included in the
independently generated IELR comparison below.

\hypertarget{62-completed-and-minimized-state-counts}{%
\subsection{6.2 Completed and minimized state
counts}\label{62-completed-and-minimized-state-counts}}

The second comparison uses Bison\textquotesingle s IELR and LALR modes,
again subtracting the synthetic accept state.

\begin{longtable}[]{@{}lrrr@{}}
\toprule\noalign{}
bare grammar & Gazelle final & Bison IELR − \texttt{\$accept} & Bison
LALR − \texttt{\$accept} \\
\midrule\noalign{}
\endhead
\bottomrule\noalign{}
\endlastfoot
C11 & 470 & 470 & 470 \\
Python & 418 & 418 & 418 \\
jq & 311 & 311 & 311 \\
SQLite & 1,340 & 1,340 & 1,319 \\
regex & 44 & 44 & 44 \\
meta & 61 & 61 & 61 \\
\end{longtable}

Five evaluated grammars require no states beyond LALR. SQLite is the
interesting case: IELR retains 21 more states than LALR to preserve
canonical resolved behavior, and Gazelle retains the same number.
Gazelle compresses the canonical machine by a factor of 49.7 while
avoiding the merges that would alter the reference parser.

The production Gazelle grammars retain their terminal modifiers. Virtual
reduce symbols and the guard of §5 can keep states separate when merging
them would manufacture a static or runtime choice absent from the
canonical machine:

\begin{longtable}[]{@{}lr@{}}
\toprule\noalign{}
production grammar & Gazelle states with modifiers \\
\midrule\noalign{}
\endhead
\bottomrule\noalign{}
\endlastfoot
C11 & 506 \\
Python & 418 \\
regex & 44 \\
meta & 61 \\
\end{longtable}

The 36-state C11 difference is therefore not a failure to reach the bare
IELR count. It is the measured cost of semantics removed from the Bison
comparison. Conversely, the equal Python, regex, and meta counts show
that modifiers do not necessarily prevent the same quotient.

The bare-count equality is empirical. It does not establish that Gazelle
and IELR always produce the same quotient, that the machines are
isomorphic, that their encoded tables occupy the same number of bytes,
or that either count is globally optimal. The present evidence supports
two separate claims: the construction preserves canonical resolved
behavior by its ordering and equivalence criterion, and its state-count
compactness matches IELR on these bare grammars.

\hypertarget{63-implementation-size-and-cost}{%
\subsection{6.3 Implementation size and
cost}\label{63-implementation-size-and-cost}}

The generic automaton module contains the NFA and DFA representations,
subset construction, iterative partition refinement, and column
equivalence used by both lexers and parsers. It is under 300 lines in
the current implementation. LR-specific code constructs the item NFA,
classifies conflicts, completes reduce edges, and extracts tables.

Gazelle eagerly allocates one NFA node for every
\texttt{(production,\ dot,\ lookahead)} triple, including unreachable
triples. This trades memory for simple index arithmetic; subset
construction visits only reachable sets. In an optimized test build, the
SQLite fixture\textquotesingle s complete Gazelle construction and Bison
IELR comparison finishes in approximately six seconds. Gazelle
constructs 66,539 canonical states and minimizes them to 1,340 in that
run; Bison\textquotesingle s canonical mode was stopped after seven
minutes without a result. The Bison measurements were reproduced with
GNU Bison 3.8.2. These figures are engineering observations rather than
a controlled performance study. A complete evaluation should report the
source revision and grammar hashes, hardware, peak memory, construction
time by phase, action and goto entry counts, generated-table bytes, and
parse speed against other generators.

The repository\textquotesingle s opt-in Bison regression covers C11,
Python, jq, regex, meta, and a small LR(1)-but-not-LALR witness. SQLite
is an extended opt-in fixture so normal development does not pay its
construction cost. An IELR-only mode runs compactness checks without
first requiring Bison\textquotesingle s potentially much slower
canonical construction.

\hypertarget{64-reproducing-the-tests}{%
\subsection{6.4 Reproducing the tests}\label{64-reproducing-the-tests}}

The ordinary Rust suite has no external parser-generator dependency:

\begin{Shaded}
\begin{Highlighting}[]
\ExtensionTok{cargo}\NormalTok{ test }\AttributeTok{{-}{-}features}\NormalTok{ codegen}
\end{Highlighting}
\end{Shaded}

With GNU Bison installed, the main differential regression runs
canonical state and conflict counts, canonical bisimulation, and final
IELR counts. A release build is recommended for jq:

\begin{Shaded}
\begin{Highlighting}[]
\VariableTok{GAZELLE\_BISON\_REGRESSION}\OperatorTok{=}\NormalTok{1 }\ExtensionTok{cargo}\NormalTok{ test }\AttributeTok{{-}{-}release} \AttributeTok{{-}{-}features}\NormalTok{ codegen }\DataTypeTok{\textbackslash{}}
\NormalTok{  bison\_canonical\_and\_ielr\_state\_counts }\AttributeTok{{-}{-}} \AttributeTok{{-}{-}nocapture}
\end{Highlighting}
\end{Shaded}

The SQLite stress fixture deliberately skips Bison canonical
construction and runs the independently generated IELR comparison:

\begin{Shaded}
\begin{Highlighting}[]
\VariableTok{GAZELLE\_BISON\_REGRESSION}\OperatorTok{=}\NormalTok{1 }\VariableTok{GAZELLE\_BISON\_EXTENDED}\OperatorTok{=}\NormalTok{1 }\DataTypeTok{\textbackslash{}}
\VariableTok{GAZELLE\_BISON\_IELR\_ONLY}\OperatorTok{=}\NormalTok{1 }\VariableTok{GAZELLE\_BISON\_GRAMMAR}\OperatorTok{=}\NormalTok{sqlite }\DataTypeTok{\textbackslash{}}
  \ExtensionTok{cargo}\NormalTok{ test }\AttributeTok{{-}{-}release} \AttributeTok{{-}{-}features}\NormalTok{ codegen }\DataTypeTok{\textbackslash{}}
\NormalTok{  bison\_canonical\_and\_ielr\_state\_counts }\AttributeTok{{-}{-}} \AttributeTok{{-}{-}nocapture}
\end{Highlighting}
\end{Shaded}

\texttt{GAZELLE\_BISON\_GRAMMAR=\textless{}name\textgreater{}} selects
any single fixture, and \texttt{GAZELLE\_BISON\_IELR\_ONLY=1} can be
used independently to run only the compact state-count comparison.

\hypertarget{7-related-work}{%
\section{7. Related work}\label{7-related-work}}

Knuth established canonical LR(k) parsing and the regularity of viable
prefixes {[}1{]}. The item-NFA account of closure and goto is
established in parsing texts and lecture treatments {[}6, 7{]}.
Gazelle\textquotesingle s contribution is not subset construction
itself, but placing reductions in the same labeled transition relation
so that later automaton algorithms see the complete resolved table.

DeRemer\textquotesingle s LALR construction merges states with equal
LR(0) cores {[}2{]}. Pager introduced compatibility tests that permit
many such merges without creating conflicts for LR grammars {[}3{]}.
IELR starts from the compact LALR structure and eliminates
inadequacies---places where merging would change the behavior of the
canonical parser after conflict resolution {[}4{]}. Gazelle targets the
same behavioral baseline from the other direction: materialize and
resolve the canonical machine, then quotient its completed behavior.

Kannapinn gives the closest prior construction {[}9{]}. He builds
completed canonical LR(k) machines and applies partition refinement
using right-context and reduction information as Moore-style state
output. He also discusses an equivalent Mealy formulation. Crucially, he
observes that ordinary minimization of the bare canonical transition
graph is insufficient because reduction information is absent from that
graph. Per-production reduce states remove precisely this obstacle.
Kannapinn assumes conflict-free LR(k) grammars; canonical conflict
resolution before minimization, conservative completion, and runtime
precedence are not part of that construction.

Reduction-incorporated generalized LR parsing also moves reductions into
an automaton, but with a different representation and purpose. Scott and
Johnstone connect reducing states to goto targets to reduce stack
activity in generalized parsing {[}10{]}. Gazelle instead uses
lookahead-labeled edges into rule-distinguished sink states so
reductions participate in table minimization.
Heilbrunner\textquotesingle s parsing-automata treatment belongs to the
broader automata-theoretic LR tradition but retains reduction
information outside the bare transition relation {[}8{]}.

langcc is a contemporary LR-family generator that attacks table size
from the construction side {[}13{]}: a continuation-passing-style
grammar transformation enlarges the accepted grammar class, and an
optimized NFA construction partitions follow sets as the automaton is
built, keeping it small rather than shrinking it afterwards. That route
and the post-hoc one taken here --- materialize, resolve, complete,
quotient --- are complementary rather than competing.

Table packing is a separate objective. Joliat studies reduced matrix
representations for LR parser tables {[}11{]}, which concerns compactly
encoding a chosen table. Default reductions are the important bridge:
the original yacc report documents per-state defaults {[}14{]}, and
Bison documents both the option and its delayed error detection today
{[}12{]}. A default is not merely an encoding of the unchanged machine;
it replaces error behavior by reductions so the resulting row can be
stored sparsely. Gazelle reuses that behavioral relaxation with an
alignment-driven selection policy, before minimization, to make states
equivalent rather than merely to make individual rows sparse. It
subsequently applies the classical per-state policy to the minimized
machine for row storage. Thus the transformation is directly related,
while the optimization objective is different.

Yang proves that minimizing LR(1) state machines is NP-hard {[}5{]}. As
§4.5 explains, Gazelle computes the coarsest behavioral quotient of one
deterministically completed machine; it does not search over
conflict-free merges or completions, and therefore does not contradict
that result --- it answers a different, tractable question whose
empirical adequacy §6 measures.

\hypertarget{8-limitations-and-future-work}{%
\section{8. Limitations and future
work}\label{8-limitations-and-future-work}}

The present construction and evaluation have six important limitations.

First, the preservation argument should be made fully formal. §4.4
reduces the parser-specific content to a single lemma --- inserted
reductions preserve non-viability of the lookahead --- with everything
after it inherited from automata theory; a proof over parser
configurations would make that lemma\textquotesingle s stack assumptions
explicit, and should include the formal statement that classical
per-state default reductions satisfy the same insertion conditions.

Second, the canonical comparison now establishes resolved-table
bisimulation, but the final comparison establishes only state-count
agreement, not final machine equivalence or equal encoded size.
Exporting Bison\textquotesingle s IELR machine into a normalization that
accounts for Gazelle\textquotesingle s completion policy would permit a
product-machine check of the final tables. Reporting action/goto entries
and serialized bytes would test the motivating claim about tables being
small enough to ship.

Third, only six grammars are measured. The two external fixtures are
mechanically normalized parser grammars, and SQLite\textquotesingle s
complemented token sets are represented by explicit wildcard-class
terminals. A corpus covering more LR(1)-but-not-LALR grammars, conflict
policies, nullable cycles, and large generated languages would better
characterize when Gazelle\textquotesingle s fixed completion matches
IELR size and when it does not.

Fourth, spurious reductions weaken immediate error detection and can
execute semantic actions on invalid input. The generated parser
preserves recognition and valid parses, not an identical failed-parse
trace. Error recovery may also observe differences that simple
recognition does not. These behaviors should be measured explicitly.

Fifth, the modifier-stripped Bison comparison isolates the minimization
algorithm but does not measure the cost of Gazelle\textquotesingle s
full resolution model. The production counts in §6.2 begin to expose
that cost; trace-equivalence tests over adversarial grammars should
cover \texttt{shift}, \texttt{reduce}, \texttt{prec}, and
\texttt{conflict} policies directly.

Finally, the current partition-refinement implementation is a simple
Moore-style fixed-point algorithm, not Hopcroft\textquotesingle s
asymptotically faster algorithm despite its function name. This does not
affect the quotient it computes, but the implementation and terminology
should be brought into agreement.

\hypertarget{9-conclusion}{%
\section{9. Conclusion}\label{9-conclusion}}

The hard part of minimizing an LR state machine is not merely deciding
which states look similar. It is preserving the meaning of conflicts
while information is discarded. Merge-first constructions must predict
or repair that interaction. Gazelle avoids it by changing both the order
of operations and the representation on which they operate. Packing the
resulting table compactly in memory is an important but orthogonal
engineering step.

Encoding \texttt{reduce\ r\ on\ a} as an edge labeled \texttt{a} into
reduce state \texttt{R\_r} makes the complete parse table a labeled
transition system. Subset construction produces canonical LR(1) item
sets and actions together. Conflicts are reachable target states and are
resolved while their canonical contexts are still intact. Conservative
completion exposes safe equivalences, and generic partition refinement
computes the coarsest quotient of the completed resolved machine.

The result is not a solution to globally minimal LR table construction.
It is a simpler construction of a canonical-LR-equivalent parser whose
state counts, on all six modifier-stripped grammars evaluated here,
equal those produced by IELR. Production grammars with additional
resolution semantics may retain more states, and equal state counts do
not imply equal encoded bytes. More broadly, the construction
illustrates a useful engineering principle: when a domain-specific
object nearly fits a mature generic algorithm, changing the
representation may be more effective than inventing another
domain-specific algorithm.

\hypertarget{acknowledgments-and-ai-assistance}{%
\section{Acknowledgments and AI
assistance}\label{acknowledgments-and-ai-assistance}}

This work originated during the ZeroMatter agentic coding hackathon.
Generative-AI systems, including OpenAI Codex and Anthropic Claude,
assisted with the implementation and review of Gazelle, the drafting and
editing of this manuscript, and literature discovery and synthesis.
Their use included generating and revising code and tests, proposing
prose and structural edits, locating candidate prior work, and helping
compare the paper\textquotesingle s claims with available documentation
and sources. The human author selected the research questions, made the
final technical and editorial decisions, ran the reported experiments,
reviewed the resulting code and manuscript, and takes responsibility for
the accuracy and integrity of the work. The AI systems are not authors.

\hypertarget{references}{%
\section{References}\label{references}}

{[}1{]} D. E. Knuth, ``On the translation of languages from left to
right,'' \emph{Information and Control} 8(6), pp. 607--639, 1965.

{[}2{]} F. DeRemer, \emph{Practical Translators for LR(k) Languages},
PhD thesis, MIT, 1969.

{[}3{]} D. Pager, ``A practical general method for constructing LR(k)
parsers,'' \emph{Acta Informatica} 7, pp. 249--268, 1977.

{[}4{]} J. E. Denny and B. A. Malloy, ``The IELR(1) algorithm for
generating minimal LR(1) parser tables for non-LR(1) grammars with
conflict resolution,'' \emph{Science of Computer Programming} 75(11),
pp. 943--979, 2010. \url{https://doi.org/10.1016/j.scico.2009.08.001}.

{[}5{]} Wuu Yang, ``Minimizing LR(1) state machines is NP-hard,''
arXiv:2110.00776, 2021. \url{https://doi.org/10.48550/arXiv.2110.00776}.

{[}6{]} D. Grune and C. J. H. Jacobs, \emph{Parsing Techniques: A
Practical Guide}, 2nd ed., Springer, 2008.

{[}7{]} J. Gallier, ``A Survey of LR-Parsing Methods: The Graph Method
for Computing Fixed Points; Computation of FIRST, FOLLOW, and LALR(1)
Lookahead Sets,'' University of Pennsylvania, 2008, §11.

{[}8{]} S. Heilbrunner, ``A parsing automata approach to LR theory,''
\emph{Theoretical Computer Science} 15, pp. 117--157, 1981.

{[}9{]} S. Kannapinn, \emph{Eine Rekonstruktion der LR-Theorie zur
Elimination von Redundanz mit Anwendung auf den Bau von ELR-Parsern},
Dissertation, Technische Universität Berlin, 2001, especially ch. 4, pp.
39--54. doi:10.14279/depositonce-276.

{[}10{]} E. Scott and A. Johnstone, ``Generalized bottom up parsers with
reduced stack activity,'' \emph{The Computer Journal} 48(5), pp.
565--587, 2005.

{[}11{]} M. L. Joliat, ``On the reduced matrix representation of LR(k)
parser tables,'' Technical Report CSRG-28, University of Toronto, 1973.

{[}12{]} Free Software Foundation, \emph{Bison: The Yacc-compatible
Parser Generator}, version 3.8.2, manual §5.8.1 ``LR Table
Construction,'' §5.8.2 ``Default Reductions,'' §5.8.3 ``LAC,'' and §5.9
``Generalized LR (GLR) Parsing'' (\texttt{\%define\ lr.type},
\texttt{\%define\ lr.default-reduction}, \texttt{\%define\ parse.lac},
\texttt{\%dprec}).

{[}13{]} J. Zimmerman, ``Practical LR parser generation,''
arXiv:2209.08383, 2022.

{[}14{]} S. C. Johnson, ``Yacc: Yet Another Compiler-Compiler,''
Computing Science Technical Report 32, Bell Laboratories, Murray Hill,
NJ, 1975.

{[}15{]} E. W. Dijkstra, \emph{ALGOL-60 Translation}, Stichting
Mathematisch Centrum, Rekenafdeling, ALGOL Bulletin Supplement nr. 10,
1961.

\hypertarget{appendix-a-survey-precedence-support-in-parser-generators}{%
\section{Appendix A. Survey: precedence support in parser
generators}\label{appendix-a-survey-precedence-support-in-parser-generators}}

This appendix records the documentation survey behind
§5\textquotesingle s comparison. For each tool, the official manual or
reference was checked (July 2026) for a mechanism that keeps both
actions of a shift/reduce conflict in deterministic LR tables and
decides between them at parse time from input-dependent data.
Generation-time precedence declarations (\texttt{\%left},
\texttt{\%right}, \texttt{\%nonassoc}, \texttt{\%prec}, and equivalents)
do not qualify: they fix each conflicted cell before the parser runs.
Mechanisms that fork the parse and rank the survivors are recorded
separately below; they decide at parse time but not on a single
deterministic stack.

\hypertarget{a1-deterministic-lr-generators}{%
\subsection{A.1 Deterministic LR
generators}\label{a1-deterministic-lr-generators}}

\begin{longtable}[]{@{}p{0.18\linewidth}p{0.19\linewidth}p{0.41\linewidth}p{0.13\linewidth}@{}}
\toprule\noalign{}
generator & construction & precedence mechanism & decided at \\
\midrule\noalign{}
\endhead
\bottomrule\noalign{}
\endlastfoot
GNU Bison (LR modes) & LALR, IELR, canonical LR(1) &
\texttt{\%left}/\texttt{\%right}/\texttt{\%nonassoc}/\texttt{\%prec} &
generation \\
Berkeley Yacc & LALR & yacc declarations & generation \\
Menhir & LR(1), Pager\textquotesingle s method &
\texttt{\%left}/\texttt{\%right}/\texttt{\%nonassoc} & generation \\
Happy (LR mode) & LALR &
\texttt{\%left}/\texttt{\%right}/\texttt{\%nonassoc} & generation \\
CUP & LALR & declarations + contextual \texttt{\%prec} & generation \\
Lemon & LALR & \texttt{\%left}/\texttt{\%right}/\texttt{\%nonassoc} &
generation \\
Hyacc & LR(1) (Pager PGM), LALR & yacc declarations & generation \\
LALRPOP & LR(1) (lane table) & none; grammar stratification & --- \\
PLY & LALR & \texttt{precedence} tuple & generation \\
SLY & LALR & \texttt{precedence} attribute & generation \\
Racc & LALR & precedence table & generation \\
Rustemo (LR mode) & LR & static declarations & generation \\
Jison (LR modes) & SLR/LALR/LR &
\texttt{\%left}/\texttt{\%right}/\texttt{\%nonassoc} & generation \\
parglare \texttt{Parser} & modified LALR (default), SLR optional & user
\texttt{dynamic\_filter} over retained actions & parse time \\
\end{longtable}

Parglare is the one deterministic precedent found. Productions marked
\texttt{dynamic} retain candidate shift and reduce actions, and a
predicate called during parsing accepts or rejects each action using the
parsing context. Its manual\textquotesingle s example makes the first
distinct operator encountered the lowest priority and later operators
progressively higher. Its filter shifts when the incoming operator ranks
higher than the pending one and reduces when it ranks lower or equal.
For example, it parses \texttt{1\ +\ 2\ *\ 3\ -\ 4\ +\ 5} as
\texttt{((1\ +\ (2\ *\ (3\ -\ 4)))\ +\ 5)}. The example instantiates the
deterministic \texttt{Parser} class. This user-supplied filter is more
general than Gazelle\textquotesingle s built-in precedence comparison,
but it establishes that runtime precedence on a deterministic LR stack
is not itself novel. Parglare\textquotesingle s deterministic tables use
a modified LALR construction by default, with SLR optional. The filter
exposes whatever actions those tables retain, and its documentation does
not relate that retained set to the canonical
automaton\textquotesingle s per-context actions, which is the
faithfulness question §5 engineers. The other thirteen surveyed tools do
not document a table entry deferring the choice at all.

\hypertarget{a2-parse-time-precedence-in-glr-systems}{%
\subsection{A.2 Parse-time precedence in GLR
systems}\label{a2-parse-time-precedence-in-glr-systems}}

GLR tables routinely retain the shift and every applicable reduction in
a conflicted cell. At runtime the parser forks over that action set, and
precedence can rank the surviving alternatives.
Tree-sitter\textquotesingle s \texttt{prec.dynamic} assigns rule weights
at runtime and selects the successful subtree with the greatest
accumulated weight. Bison\textquotesingle s GLR mode uses
\texttt{\%dprec} to prefer the surviving action with the highest
declared dynamic precedence {[}12{]}. Lezer similarly permits declared
ambiguities resolved during its GLR-style parse, and Happy has a GLR
mode with semantic disambiguation. In each case the mechanism
presupposes nondeterministic parsing; none applies to the
tool\textquotesingle s deterministic LR tables.

\hypertarget{a3-runtime-precedence-outside-lr-generation}{%
\subsection{A.3 Runtime precedence outside LR
generation}\label{a3-runtime-precedence-outside-lr-generation}}

Runtime-defined operators exist in settings that are not LR table
generation. Prolog\textquotesingle s \texttt{op/3} directive alters the
operator table of its reader, an operator-precedence parser. Pratt
parsing and precedence climbing decide from a binding-power table at
parse time, in hand-written recursive descent. Swift\textquotesingle s
\texttt{precedencegroup} and Haskell\textquotesingle s fixity
declarations feed precedence resolution phases inside those compilers.
These confirm that the \emph{language feature} is in demand. Parglare
additionally shows that a general runtime action filter can provide it
in a deterministic LR parser. Gazelle\textquotesingle s narrower
contribution is the token-carried table mechanism and its preservation
through completion and minimization.

\hypertarget{a4-a-witness-merging-defeats-table-local-filtering}{%
\subsection{A.4 A witness: merging defeats table-local
filtering}\label{a4-a-witness-merging-defeats-table-local-filtering}}

To make the faithfulness distinction concrete, we constructed a
five-production grammar and ran it against parglare 0.21.1
(\texttt{blog/parglare\_witness.py} in the repository):

\begin{Shaded}
\begin{Highlighting}[]
\NormalTok{S: \textquotesingle{}a\textquotesingle{} M | \textquotesingle{}b\textquotesingle{} M Z;}
\NormalTok{M: E \{dynamic\} | E Z \{dynamic\};}
\NormalTok{E: \textquotesingle{}x\textquotesingle{};}
\NormalTok{terminals Z: \textquotesingle{}z\textquotesingle{} \{dynamic\};}
\end{Highlighting}
\end{Shaded}

Canonical LR(1) gives the state reached by \texttt{a\ E} an
unconditional shift on \texttt{z} --- the completed
\texttt{M\ -\textgreater{}\ E} carries lookahead \texttt{\{EOF\}} there
--- while the same-core state reached by \texttt{b\ E} has a genuine
shift/reduce question on \texttt{z} (\texttt{b\ x\ z} must reduce,
\texttt{b\ x\ z\ z} must shift). parglare\textquotesingle s construction
merges the two states and unions their lookaheads, so one cell answers
for both contexts. Under the default shift preference, the question is
resolved away at construction despite the \texttt{dynamic} markers, and
the valid input \texttt{b\ x\ z} cannot be recovered by filtering the
remaining candidate. Notably, parglare\textquotesingle s
\href{https://github.com/igordejanovic/parglare/blob/0.21.1/docs/disambiguation.md\#dynamic-disambiguation-filter}{0.21.1
manual} shows its dynamic-precedence example constructed as
\texttt{Parser(grammar,\ dynamic\_filter=...)}, but the
project\textquotesingle s corresponding
\href{https://github.com/igordejanovic/parglare/blob/0.21.1/tests/func/parsing/disambiguation/test_dynamic_disambiguation_filters.py}{regression
test} adds \texttt{prefer\_shifts=False}; without it, the documented
results are not obtained. With \texttt{prefer\_shifts=False} the merged
cell defers --- in both contexts: at the decision point,
\texttt{a\ x\ z} and \texttt{b\ x\ z} present the same automaton state
and the same remaining input yet require opposite actions, so a shift
policy fails \texttt{b\ x\ z} and a reduce policy fails
\texttt{a\ x\ z}. No filter that reads only the state, candidate action,
semantic subresults, and unread input can be correct for both.

Parglare\textquotesingle s actual callback is not limited to that local
information: it can inspect the full input, retain arbitrary state in
\texttt{context.extra}, or run arbitrary code. The executable witness
therefore also includes a history-aware filter that distinguishes the
two inputs by rereading their consumed prefix and parses all four
successfully. More generally, a callback could maintain a shadow
canonical parser. That escape does not restore faithfulness to the LALR
table; it makes user code reconstruct the context that merging erased.
Counting such an unrestricted callback as a complete answer would reduce
parser generation to the oracle construction excluded above. A
canonical-faithful table instead keeps the contexts apart --- the
\texttt{a} cell a plain shift, the \texttt{b} cell the one deferred
entry --- which is what §5\textquotesingle s guard preserves through
completion and minimization. The witness concerns this grammar and
parglare 0.21.1; parglare\textquotesingle s GLR parser, which forks
instead of choosing, parses all four inputs without such reconstruction.

\hypertarget{a5-scope}{%
\subsection{A.5 Scope}\label{a5-scope}}

The survey records documented mechanisms in the tools and versions
examined; its negative entries do not establish impossibility, and an
undocumented or later mechanism would change the comparison.

The deterministic-generator rows were checked against the following
official manuals and project references: GNU Bison 3.8.2 {[}12{]};
\href{https://invisible-island.net/byacc/manpage/yacc.html}{Berkeley
Yacc\textquotesingle s manual}; the
\href{https://gallium.inria.fr/~fpottier/menhir/manual.html}{Menhir
Reference Manual}, version 20260209, §§4.1.4, 6, and 18;
Happy\textquotesingle s
\href{https://haskell-happy.readthedocs.io/en/latest/using.html\#how-to-specify-precedences}{LR
precedence documentation} and
\href{https://haskell-happy.readthedocs.io/en/2.0.2/glr.html}{GLR
documentation}; the
\href{https://www.cs.princeton.edu/~appel/modern/java/CUP/manual.html}{CUP
manual}; the \href{https://sqlite.org/lemon.html}{Lemon documentation};
the \href{https://hyacc.sourceforge.net/}{Hyacc project documentation};
the
\href{https://lalrpop.github.io/lalrpop/tutorial/004_full_expressions.html}{LALRPOP
expression tutorial}; the
\href{https://ply.readthedocs.io/en/latest/ply.html}{PLY documentation};
the \href{https://sly.readthedocs.io/en/latest/sly.html}{SLY
documentation}; the \href{https://ruby.github.io/racc/index.html}{Racc
reference}; the
\href{https://igordejanovic.net/rustemo/grammar_language.html}{Rustemo
grammar documentation}; the
\href{https://gerhobbelt.github.io/jison/docs/}{Jison documentation};
and parglare\textquotesingle s
\href{https://igordejanovic.net/parglare/stable/disambiguation/\#dynamic-disambiguation-filter}{dynamic-disambiguation
documentation}.

The comparisons outside deterministic LR use
tree-sitter\textquotesingle s
\href{https://tree-sitter.github.io/tree-sitter/creating-parsers/2-the-grammar-dsl.html}{grammar
DSL reference}, the
\href{https://lezer.codemirror.net/docs/guide/\#ambiguity}{Lezer guide},
and SWI-Prolog\textquotesingle s
\href{https://www.swi-prolog.org/pldoc/man?predicate=op/3}{\texttt{op/3}
reference}. All web sources were accessed in July 2026.

\end{document}
